\documentclass{ecnreport}

\stud{OD Robotique 2025-2026}
\topic{ROS 2 exam}

\begin{document}

\inserttitle{ROS 2 exam}

\insertsubtitle{2h, documents / internet allowed (closed chats please)}


\section{Description}

In this exam you will have to write a C++ node and run it several times through a launch file.\\

The package should first be compiled, and you can configure the IDE with:
\begin{bashcodelarge}
 cd ~/ros2/src/ecn_2025 # or wherever you have put the repo
 colbuild -t            # compile this package
 ideconf                # generate the IDE config file
\end{bashcodelarge}

Then the simulation can be run once and for all with:
\begin{bashcodelarge}
 ros2 launch ecn_2025 simulation_launch.py
\end{bashcodelarge}

This simulation involves the Baxter robot which has to perform visual servoing with each of its arms.
In order to do this, you have to publish a joint velocity command for the arm in order to align the target with the camera frame.


\section{Control node}

The control node (\okt{control\_node.cpp}) should be run for each arm and should:
\begin{itemize}
 \item Retrieve the  pose of the target with regards to the camera
 \item Retrieve the ${}^e\mathbf{J}_e$ Jacobian through a service
 \item Compute the control law and publish it as joint velocity
\end{itemize}

All examples below use the left arm, just replace with \okt{right} for the other limb.

\subsection{Relative pose between target and camera}

The frames are named as follows:
 \begin{itemize}
  \item the camera frame name is \okt{left\_hand\_camera}
  \item the target frame name is \okt{left\_target}
  \end{itemize}

  Use the provided TF \okt{buffer} in order to check is the transform is available and to get it if it is the case.
  Note that we are only interested in the \okt{Transform}, while \okt{buffer.lookupTransform} returns a \okt{TransformStamped}. Make sure that the transform you get is in the correct direction (e.g. Z should be positive as the object is in front of the camera).

\subsection{Calling the Jacobian service}

The simulation proposes a service for each arm to retrieve the current Jacobian. It is of type \okt{baxter_simple_sim::srv::Jacobian}. Use the following command to have an overview of this service:
\begin{bashcodelarge}
ros2 interface show baxter_simple_sim/srv/Jacobian
\end{bashcodelarge}

In our case we want to set the \okt{ee_frame} field to \okt{True} and do not provide any joint position, as we use the current ones.

The service can be called at the following address: \okt{/robot/limb/left/jacobian}.\\
As in the lab 3, we use a wrapper (namely \okt{ServiceNodeSync}) to have synchronous service calls. This wrapper should:
\begin{itemize}
 \item be initialized with a unique name and the address of the service though the \okt{init} function
 \item call the service with the \okt{call} function that returns a \okt{std::optional<Response>}.
\end{itemize}
If the returned \okt{std::optional} has a value (\okt{response.has_value()}) then the Jacobian is accessed with \okt{response->jacobian}.

\subsection{Visual servoing control law}

The function \okt{computeVisualServo} takes the pose of the target in the camera frame and the end-effector Jacobian. It returns a \okt{std::vector<double>} for the joint velocity that tracks the object. This value should be written in the command message and published on the well-known topic \okt{/robot/limb/left/joint_command}.

\section{The launch file}

Once your node runs for a given arm, add a parameter to pick the side (left or right) and run it twice through the \okt{both_launch.py} launch file.

Give two separate names to the node or there might be a name conflict. You do not need to use any namespace, the side parameter should be enough to handle all topics, services and frame names.


\end{document}
